\label{sec:introduction}

The 2020 redistricting cycle demonstrated that census geographic units are not
created equal for redistricting purposes. Census tracts---the primary unit for
\textsc{bisect} redistricting---average approximately 4{,}000 people with a typical
area of a few square miles. They are designed as statistical units, not as building
blocks for congressional districts. For states with many congressional seats, the
tract resolution is adequate: California ($k=52$, $n=9{,}769$ tracts, ratio $= 0.0053$)
has nearly 200 tracts per congressional district on average, giving the bisection
algorithm abundant flexibility to achieve population balance. For states with few
congressional seats, the tract resolution creates problems: Wyoming ($k=1$, $n=157$
tracts, ratio $= 0.0064$) is an extreme case but illustrates the general pattern
that small-delegation states have fewer tracts per district.

The F.3 paper \citep{f3paper} established the $k/n > 0.05$ resolution rule through
empirical analysis of redistricting failures in high-$k/n$ state legislative chambers.
When this ratio exceeds 0.05, the recursive bisection tree encounters partitions where
the two sub-graphs have incompatible population levels---the algorithm must split
a single tract between two districts, violating the geographic integrity requirement.
The resolution rule says: upgrade to block groups when this happens.

For congressional redistricting in 2020, 39 states required block-group resolution.
In 2030, the number rises to 43. This paper documents the technical implications
of that expansion.

\subsection{Block Groups vs.\ Tracts: Geographic Definitions}

Census block groups are a geographic level between tracts and individual blocks.
By definition, every block group is entirely contained within a single tract,
and every tract contains between one and nine block groups (median: four).
Block groups average approximately 1{,}000--3{,}000 people and are designed to
match recognizable neighborhood boundaries in urban areas.

For redistricting purposes, the key advantage of block groups over tracts is their
smaller size: a congressional district that contains only 15 tracts may contain
60 block groups, providing the bisection algorithm with 4$\times$ more partition
choices. The key disadvantage is data volume: 240{,}000 tracts nationally vs.\
roughly 1.2 million block groups. Adjacency graphs at block-group resolution are
5$\times$ larger, requiring proportionally more computation and storage.

\subsection{Relationship to Q-Series}

This paper is Q.3 in the Forward-2030 series. It assumes the reapportionment
projections from Q.4 (which provide the $k$-values used in the resolution rule
application) and the data update pipeline from Q.1 (which provides the ACS and
LODES data at block-group resolution). The construction timeline is coordinated
with the pre-registration schedule documented in Q.0 (overview) and Q.2
(algorithm pre-registration).
